\section{Conclusion}
\label{sec:conclusion}

\subsection{Key findings}
% TODO
\subsection{Limitations}
The shortcomings of the work are the extractor's weakness in removing interfering speech, whilst simultaneously not introducing artefacts that get misinterpreted by the downstream model as target speech.

\subsection{Future work}
The structure and nature of the approach formulated in this work allows for a number of future directions to be explored.   
\begin{itemize} 
      \item \textbf{Streaming architecture.} The LSTM layers of the extractor could be replaced with state space layers. A state space layer is a recurrent layer whose updates are linear, which allows for parallel training while still running frame by frame. SA-Mamba \cite{liu2026samamba} applies this idea within a band-split, dual-path design similar to the one in this work, and conditions the state space layers on the target speaker. It placed second in the online track of the REAL-TSE challenge, while its report describes training on simulated mixtures only, and the
  organisers attribute most of the gains of the top systems to their data \cite{wang2026realtse}. The architecture has not yet been used with a live speech model as the listener and secondly it did not have a training recipe built around one, as in this work. 
      \item \textbf{Training objective.} The current objective is a proxy for general extraction rather than for the downstream task. Weighting the leaked speech or the artefacts in the loss moved errors from one source to the other but did not lower the LCF-WER (\autoref{sec:results_metric}), so an objective closer to \emph{how} Gemini listens is still needed.
      \item \textbf{Removing interfering speech without artefacts.} As seen in the results, the extractor struggles with the trade-off between removing interfering speech and introducing artefacts. A likely reason is the aggressive masking, which can leave sharp transitions where the interfering speech was removed. The extractor could instead be trained to also extract the interfering speaker, so that it learns what to remove as well as what to keep.
      \item \textbf{Filling the gaps for Gemini.} Rather than leaving gaps where the interfering speech was removed, the extractor could fill them with a signal that Gemini ignores, such as low-level noise, leaving nothing for Gemini to fill with words. Gemini, during experimentation, was found to be robust to very low-level noise, but the idea is to make audio clips that are optimised to be ignored by the listener.
      \item \textbf{Adaption to real recordings.} Although fine-tuning to the live transcriber did not yield improvements, the extractor could be adapted to real recordings after training on simulated mixtures.
      \item \textbf{Learning the band plan.} No paper has yet justified the actual boundaries of the bands used in the band-split design. Learnable band positionings could find better boundaries for the extractor and could be conditioned on the target speaker. This will introduce another conditioning to the extractor, which by the extension in this work showed to be beneficial.
  \end{itemize}
